\uuid{RTh7}
\titre{Noyau/Image/Injection/Surjection/Théorème du rang}
\auteur{hanani ; mzari}
\organisation{exo7}
\licence{CC BY-NC-SA 4.0}
\datecreate{2020-01-20}
\niveau{L1}
\module{Algèbre}
\chapitre{Application linéaire}
\sousChapitre{Définition}
\difficulte{3}
\theme{Noyau et image, Noyau, Image, Injection, Surjection, Théorème du rang}
\exo7qcmid{lille-2.applications-lineaires.024}

\contenu{
\begin{question}
\qclassification{107.01}
\qtags{motcle=Noyau/Image/Injection/Surjection/Théorème du rang}
Soit $E$ et $F$ deux $\Rr$-espaces vectoriels de  dimensions finies tels que $\dim E= \dim F=n$ et $f$ une application linéaire de $E$ dans $F$. Quelles sont les assertions vraies ?
\begin{answers}
\bad{Si $\dim \ker f =0$, alors $\dim \Im f <n$}.
\good{si $f$ est injective, alors $f$ est surjective}.
\good{Si $\dim \Im f <n$, alors $\dim \ker f >0$}.
\good{si $f$ est surjective, alors $f$ est injective}.
\end{answers}
\begin{explanations} Du théorème du rang, on déduit que $f$ est bijective si, et seulement si, $f$ est injective ou surjective.
\end{explanations}
\end{question}
}
